\documentclass[english,11pt,a4paper]{article}
\usepackage[T1]{fontenc}
\usepackage{babel}
\usepackage{cleveref}
\usepackage{amssymb}
\usepackage{fullpage}

\newcommand{\comma}[2]{#1\hspace{1pt} {\downarrow}#2}
\newcommand{\id}[1]{\mathsf{id}_{#1}}
\newcommand{\Set}{\mathbf{Set}}
\newcommand{\pro}{\mathsf{prod}}

\title{``EGGs are adhesive!''\\
	Answer to reviewers}
\author{R. Biondo, D. Castelnovo, F. Gadducci}


\begin{document}
	\maketitle
	
	
We wish to thank the reviewers for the time they have dedicated to our paper, and for the many useful suggestions. We corrected all typos indicated; below we explain how we have addressed the issues they raised.

We have rewritten and expanded Section 6 in response to the reviewer's request. In particular, we have made explicit which kinds of DPO rules we are interested in and which issues may arise in this setting.

Regarding Sections 3 and 4, we have chosen not to alter their structure. The order of presentation has been chosen with the following rationale: we first introduce plain hypergraphical structures, then add equivalences and finally study EGGs. We believe this provides the most straightforward way of presenting the material.

We also consider a low-level, set-theoretic presentation of most of the structures under study to be necessary. Indees, the notion of separated and pruned hypergraph, which provides a possible hypergraphical basis for the notion of EGGs, is inherently formulated in terms of elements and cannot be easily treated entirely at the categorical level.

Besides these changes, Reviewer 2 raised a serious technical issue with our previous Lemma 5.7 (now Lemma 5.9).  Indeed, the lemma was incorrect in its previous form. The issue is that there is no reason for the left face of the cube used in the proof to be a pullback. We resolved this issue by imposing an additional condition on the class $\mathsf{Pb}$ (see new Definition 5.7). Luckily, this does not affect the core of our results, now summarised 
at the beginning of Section~6.


\section*{Reviewer 1}

\begin{itemize}
	\item \emph{* "an hypergraph" -> "a hypergraph"}
	
	Fixed.
	
	\item \emph{* "cfr." -> "cf."}
	
	Fixed.
	
	\item \emph{ p. 1: "then it is unique" => clarify that here and elsewhere "unique" is mathematically speaking an imprecise formulation (i.e., it should instead read "essentially unique" or "unique up to unique isomorphism")}
	
	Fixed.
	
	
	\item \emph{p. 1: "if only the left-hand side belongs to $\mathcal{M}$, the theory is still under development ..." => this is not quite correct, as per reference [5] in your list of references (which studied the "non-$\mathcal{M}$-linear" cases in full detail, not just the "semi-$\mathcal{M}$-nonlinear" cases)}
	
	We added a reference to [3] and [5]. Notice that in [5] the matches belong to the class $\mathcal{M}$, while in this paper this is not the case (see the new sec. $6$ for an example).
	
	\item \emph{p. 2: "we put the machinery at work" -> " we put the machinery to work"}
	
	Fixed.
	
	\item \emph{ "and consequently so Sections 4.1.1, 4.2, ... " -> "and consequently so are Sections 4.1.1, 4.2, ... "}
	
	Fixed.
	
	
	\item \emph{ "called respectively the sets" -> "called, respectively, the sets"}
	
	Fixed.
	
	
	\item \emph{(cfr. Section 5.2 and ??)" => please fix the erroneous cross-reference "??"}
	
	Fixed.
	
	\item \emph{We define two functions $in_H, out_H : V_H \times E_H \Rightarrow \mathbb{B}$" -> "$\mathbb{B}$" appears to be a typo (as you used the notation "$\mathbb{N}$" before)}
	
	Fixed.
	
	
	
	\item \emph{ "is left-monogamous but not separated" => you have not yet defined the notion of a hypergraph being "separated" at this point}
	
	We moved the example after  def $3.37$.
	
	\item \emph{ p. 8, Example 3.25: you have multiple notational errors in this near trivial example, which are confusing to the attentive readers: (1) you speak of "hypergraphs G and H below" but seem to refer to the diagram set to the right of the paragraph, where hypergraphs are not named; (2) "$out_G(a, e) = 2$" -> "$out_G(v, e) = 2$" (there is no node called "a" in the diagram), (3) for the hypergraph $H$, you speak of "$out_H(u, e_1)=out_H(u, e_2)=1$", yet as per your notational conventions, only right-pointing links between hyperedge boxes and nodes count for the out-degree function (i.e., the way you drew the diagram for $H$ you would need to read off "$out_H(u, e_1)$", "$out_H(u, e_2)=0$", "$in_H(u,e_1)=0$", and "$in_H(u,e_2)=1"$}
	
	We fixed the notational errors and changed the graphical notation adding an arrowhead to denote the targets.
	
	\item \emph{ p. 8: "The remaining of this section" -> "The remainder of this section"}
	
	Fixed.
	
	\item \emph{ p. 8: "Being $LMH$ a full subcategory of $Hyp$," -> "Since $LMH$ is a full subcategory of $Hyp$, "}
	
	Fixed.
	
	\item \emph{ p. 8: "Definition 3.31. A morphism $(f, g) : G \to H$ of hypergraphs is downward-closed if $t_H(h)$ has a letter belonging to the image of $g$, then $h$ is in the image of $f$." => this is rather ill-formulated - please clarify this definition (you presumably intended to express that a downwards-closed morphism is one such that, whenever a letter $l$ occurring in $t_H(h)$ for some hyperedge $h$ in $H$ is in the image of $g$, then the hyperedge $h$ itself must be in the image of $f$; coincidentally, perhaps it would be best to use different notations for morphisms of hypergraphs and elements of sets of hyperedges, since clearly $f,g,h$ notations does not make sufficiently explicit that $f$ is the mapping of hyperedges, $g$ the mapping of nodes, and $h$ a hyperedge in $H$)}
	
	We rewrote the definition.
	
	
	
	\item \emph{ p. 9, Proposition 3.33: there is something odd notationally about the occurrences of a seemingly spurious symbol "M", as in "If $M(f_1, g_1)$ is mono and downward-closed", which by the preceding examples of notations should read "If $(f_1, g_1)$ is mono and downward-closed" (i.e., no "$M$" occurring); either in some previous version of these materials you had a definition and motivation for some operation $M$, or this is a systematic) typo => please fix this carefully!}
	
	Fixed.
	
	\item \emph{ p. 9: "We refer again the reader to" -> "We refer the reader to"}
	
	Fixed.
	
	\item \emph{ p. 9: "Putting together ..." -> "Combining ..."}
	
	Fixed.
	
	\item \emph{ p. 9: "in which we are interested in" -> "that we are interested in"}
	
	Fixed.
	
	\item \emph{ p. 9: "amount to" -> "amounts to"}
	
	Fixed.
	
	\item \emph{ p. 9: "disjoint from, respectively ..." -> "disjoint from, respectively, ..."}
	
	Fixed.
	
	\item \emph{ p. 9: "which "split" a given hypergraph" -> "which "splits" a given hypergraph"}
	
	Fixed.
	
	\item \emph{ p. 10: "collects all the operation" -> "collects all the operations"}
	
	Fixed.
	
	\item \emph{ p. 11: "we conclude our exam" -> "we conclude our study"}
	
	Fixed.
	
	\item \emph{ p. 11: "while will call the third one" -> "and call the third one"}
	
	Fixed.
	
	\item \emph{ p. 11: "If we restrict to pruned separated hypergraph" -> "If we re
		strict to pruned separated hypergraphs"}
	
	
	Fixed.
	
	\item \emph{p. 11: "Lemma 3.54. $pSep$ is equivalent to $Set/(\mathbb{N}\times \mathbb{N})$." => you never clarified explicitly in your manuscript in which concrete mathematical sense you consider $\mathbb{N}$ a category, even though the proof in the appendix provides some hints at what kind of construction you mean to denote; there are also some related imprecisions in your proof, so this result must be more carefully and more precisely presented}
	
	$\mathbb{N}$ is just a set, and $Set/(\mathbb{N}\times \mathbb{N})$ is the comma category over the set $\mathbb{N}\times \mathbb{N}$.
	
	\item \emph{ p. 13, Remark 3.69 => This kind of overview of materials is presented much too late in the manuscript, and should, in fact, be the main introduction to the technical materials section.}
	
	We moved the remark at the beginning of the section, with some explanatory remarks.
	
	
	\item \emph{ p. 16: "The category $Hyp_{eq,G}$ of hypergraph" -> "The category $Hyp_{eq,G}$ of hypergraphs"}
	
	Fixed.
	
	\item \emph{ p. 16: "pruned ans separated" -> "pruned and separated"}
	
	Fixed.
	
	
	\item \emph{ p. 18: "we have that such class" -> "we have that this class"}
	
	Fixed.
	
	
	\item \emph{ p. 19: "pushouts alongs arrows" -> "pushouts along arrows"}
	
	Fixed.
	
	
	\item \emph{ p. 20: "to be simulates" -> "to be simulated"}
	
	Fixed.
\end{itemize}


\section*{Reviewer 2}

\begin{itemize}
	
	\item \emph {p.5, Lem.2.12. Is there a unique arrow $k_h$ making all three
		diagrams commute or is there a unique arrow for each diagram
		separately?
		$h\times h$ goes from $X\times X$ to $Z \times Z$, this has to be fixed.}
	
	We fixed the typing of $h \times h$. The statement of the lemma is that there exists a unique arrow making all three diagrams commute. Notice, however, that commutativity of the third diagram is equivalent to that of the other two.
	
	\item \emph{ p.6, Prop.3.1: What are the $f^n$?
		I think that the introduction of the star operator is a bit overly
		complicated. Do you lose anything if you just define $X^\star$
		via the standard Kleene star (all finite words over $X$)?}
	
	The explicit construction of the Kleene star is precisely the formalization of  ``(all finite words over $X$''). We made explicit what $f^n$ are.
	
	\item \emph{p.7, Prop.3.7: As far as I can parse the definition the 
		objects ``$\comma{\id{\Set}}{\pro^\star}$''  are maps $X\rightarrow X^\star \times X^\star$.
		But are these really hypergraphs? Shouldn't you distinguish
		between nodes and edges? Where is the proof of this proposition?}
	
	This is considered kind of folklore. An hypergraph is a pair of arrows $E \to V^\star$, or equivalently, an arrow $E \to V^\star \times V^\star$, that is an object of the comma category. We have added a reference.
	
	
	\item \emph{ p.8, Lem.3.17 is a bit hard to parse and digest. Can you make
		it more accessible?}
	
	We agree on the difficulties, but at the same time the lemma has a rather technical nature, and we preferred to preserve its current formulation. 
	
	\item \emph{ p.8, Def.3.21: What is $\mathbb{B}?$}
	
	Fixed.
	
	
	\item \emph{ p.9: You should make a graphical distinction of the source
		and target nodes of a hyperedge.}
	
	Fixed.
	
	\item \emph{ p.10, Rem.3.39: doesn't the third item subsume the second one?}
	
	No, the second one exclude the presence of cycles, the third one assure us that every edge is not connected to other edges. 
	
	\item \emph{ p.11, Rem.3.42 and Rem.3.49: In my opinion theses remark, as well as
		this entire section of the paper, is somewhat too detailed, spelling
		out relative obvious observations. It would be better to remove them
		or move them to the appendix and instead spend more time on
		e-hypergraphs and rewriting.}
	
	We moved Rem 3.49 to the appendix (now is Rem. A.6). However, we consider Rem 3.42 helpful for the intuition and we left it in the main body.
	
	\item \emph{ p.12: Sep and pSep are rather degenerated categories. It would be
		helpful for the reader if you could explain in more detail what you
		are doing here and what the motivation for considering those categories
		is.}
	
	The idea here is to provide different hypergraphical bases for EGGs. Since we have equivalences at our disposal, hyperedges need not actually share a node.
	
	\item \emph{ p.12. Prop.3.50: What do $Pr$ and $sp_P$ do? (pr=pruning?)}
		
		Indeed pr stands for pruning. The definition of the two functors is contained in the proof of $3.49$.
		
		\item \emph{ p.13, Rem.3.60: What is $U_{Hyp}$?}
		
		$U_{Hyp}$ is the forgetful functor sending an hypergraph to is set of nodes. It is defined in  Proposition 3.10
		
		\item \emph{ p.14: The addition of equivalence - a central concept - comes fairly
			late in the paper. As said above, I believe that the sections
			before can be shortened. The same holds to some extent for the
			results on hypergraphs with equivalences, which are somehow just
			hypergraphs with an explicit surjective labelling. It is e-graphs
			that are actually interesting.}
		
		Technically, the preceding sections contain results which are needed for that section, and we were at loss in identifying possible
		shortenings. We also believe that the kind of story we are presenting, as now summed up at the beginning of Section 6, may be 
		of some value for the community working of Adhesive rewriting syetems,
		
		\item \emph{ p.15: I would give concrete constructions for the pushouts in Hyp.}
		
		Pushouts for Hyp are constructed componentwise, as for Lemma B.3.
		
		\item \emph{ p.18, Prop.5.2: At this point kernel pairs are finally used (after
			having been introduced long before). It would be good to explain
			their role on a more intuitive level.}
		
		Kernel pairs are nowadays a standard tool in the categorical study of relations, thus we prefer to introduce them in Sec. 2, alongside the other general categorical preliminaries.
		
		\item \emph{ p.21: The last paragraph of Sct.6.1 is a bit cryptic and should be
			explained in more detail. This should actually be the motivation for
			considering the various variants of hypergraphs, so this should be
			carefully explained to the reader.}
		
		We rewrote Sec. 6.
		
		\item \emph{ p.21 ff.: As stated above, it is not so clear to me which rules
			you are actually interested in. Do you want rules that can delete
			items or only rule that add something?}
		
		We rewrote Sec. 6.
		
		\item \emph{ p. 23: The comparison to [21] (a paper by Ghica, Barrett and Tiurin)
			is a bit weak. If [21] is in some sense more general than your
			paper, there should be a way to make a deeper comparison.}
		
		We rewrote the paragraph comparing the two works.
		
	\end{itemize}


	
\end{document}